\chapter{Architecture of a Type-Class Solver in \rocqelpi}

Type classes provide an elegant mechanism for implementing ad-hoc
polymorphism~\cite{wadler1989}. They allow the same symbol to be reused in
different contexts while preserving the same mathematical properties. From a
Prolog-like perspective, a type-class problem corresponds to resolving a query
to a predicate that associates its input with an instance solving the
type-class constraint. In \rocq, instances act as rules stored in a database,
which are used to backchain and resolve type-class queries.

To address type-class unification problems, we extended \rocq with an API called
\cI{Class_tactics.register_solver}. This API allows developers to register
custom type-class solvers within the \rocq elaborator. In this way, type-class
resolution can be delegated to a meta-language: the custom solver computes a
solution (if it exists) to the type-class constraint and returns it to the
\rocq elaborator. We refer to such a custom solver as a \emph{meta-solver}.

% When the elaborator encounters a type-class constraint during the elaboration
% of a \rocq term, it may delegate its resolution to a registered meta-solver.
% The meta-solver receives the type-class goal together with the context in
% which it appears and must return an updated environment (called an
% \cI{evar_map} in \rocq terminology).

To support this mechanism, we also introduced another API,
\cI{Classes.register_observer}. This API provides a systematic way to notify
the meta-language whenever a new class or instance is declared in \rocq. The
meta-language can then translate these declarations into its own
representation and store them in its knowledge base.

Concretely, a meta-language integration must provide two components:
(1) a compiler that translates \rocq class and instance declarations into the
meta-language representation, and (2) a type-class solver that animates the
search by interacting with the resulting database.

In our setting, the meta-language used to implement the meta-solver is
\elpi. Consequently, the main task is to implement the class and instance
compiler, which incrementally builds the database on which the \elpi program
operates. 

The type-class solver itself typically follows a \prolog-like resolution
strategy: instances are tried according to their priority, and backtracking
may reconsider previous choices when a branch fails. In \elpi, this behavior
comes for free, since the resolution mechanism of the type-class solver
strongly imitates the standard execution model of the \elpi runtime.

In the following sections
(\cref{sec:class-encoding,sec:instance-encoding,sec:goal-encoding}), we
present the compiler by separating its three main components: the class
compiler, the instance compiler, and the goal compiler.

In \cref{sec:rule-indexing}, we describe our work on an alternative rule
indexing mechanism for the \elpi runtime, designed to speed up clause
retrieval during resolution.

Finally, in \cref{sec:tc-discussion}, we discuss practical aspects of using
the type-class solver. In particular, we explore how users can control and
customize the solver to better fit their needs. We also demonstrate the
application of our solver to existing libraries, such as \tlc, \stdpp, and
\iris.

\section{Class encoding}
\label{sec:class-encoding}

A type-class goal is a \rocq type\footnote{Or term, in \rocq terms and types are
confused} $T$ whose head is a type-class $C$. The resolution of such a goal
consists in finding an \rocq instance whose type is an instantiation of $T$. The
goal $T$ has in general $n$ arrows, one per the declared arguments of $C$.

Ideally a type-class goal should be translated in a query with an
extra-argument. This extra-argument shold carry the \rocq instance for the
initial goal.

In the \elpi compiler, we create a new predicate for each new class. This
predicate carries the name of the class and has $n+1$ arguments. The $n$ first
arguments represent the parameters of the class itself. The last extra argument
is a term carrying the concreate instance term representing the solution of the
query.

As an example, we take back the \cI{Add} type-class
introduced in \cref{sec:ad-hoc-poly}. For convenience, the full definition of
the class and its instances is reproduced in \cref{fig:add-class-full}.

\begin{figure}
  \input{code/plus_tc_add_inst.tex}
  \caption{Add class and instances}
  \label{fig:add-class-full}
\end{figure}

% \cI{Add} carries a parameter \cI{T} of type \cI{Type}. Ideally a \rocq goal
% \cI{Add T}, where $T$ is a term, should corresponds to an \elpi query with the
% following shape \eI{tc-Add {{T}} Sol} where \eI{Sol} is a unification variable
% that is instantiated by backchain the query into the database. It means that
% the compiler

% Concretely, in the \elpi
% settings, we need to create a new predicate to be asssociated to the \cI{Add}
% class. This predicate should carry on argument

When declaring \cI{Add}, the \elpi class compiler is triggered and receives
as input the \rocq symbol \cI{global (indt «Add»)}. Thanks to the 

\section{Instance encoding}
\label{sec:instance-encoding}

register hook
\subsection{From term to clause}

In a first version we could be naif and have: tc-class TERM,
but this is not good for X reasons:

\begin{itemize}
  \item We cannot benefit from modes
  \item We cannot benefit from indexing
  \item We cannot benefit from determinacy checker
\end{itemize}

\subsection{Conditional instances: encoding premises}

\subsection{Rule priorities}

\section{Goal encoding}
\label{sec:goal-encoding}

Devo compilare il contesto

\section{Rule indexing}

\subsection{Discrimination trees}
\label{sec:rule-indexing}

\newcommand{\pfdt}{https://www.cs.cmu.edu/~fp/courses/99-atp/lectures/lecture28.html}

Discrimination trees: Pfenning notes~\footnote{\url{\pfdt}}

TODO: do some benchmarks

Discrimination trees: opt for lists and tree depths.

\subsection{Some ideas on tabling}

\cite{tabling-dm,tabling-pie}


\section{Discussion}
\label{sec:tc-discussion}

\subsection{User defined rules}
Hint extern etc

\subsection{Accessing the database}

\cite{fissore2023}

\subsection{Experimental results}

\subsection{Unification problems}

\paragraph{$\eta\beta$-unification}
\paragraph{$\delta$-unification, coercions and canonical structures}

